\section{Discussion and Limitations}
\label{sec:discussion}
\textbf{What to report.} The evidence supports a conditional answer. The naive report of a searched system is biased upward whenever more than one configuration is compared on the same items; the bias is of the order of the true differences between configurations when $n$ is a few hundred and $K$ is large, so reported improvements of a few points are not interpretable without $n$ and $K$. Reporting the score of $\hatc$ on items not used for the search (M1) removes the bias at the price of items; a conditional/hybrid selection-aware interval (M4) gives nominal coverage for exhaustive search without extra items but needs the full score table and, in our runs, is wide; cross-fitting (M2) estimates the performance of the procedure, with mildly undercovering intervals. At minimum a report should state $K$, $n$ and the search procedure, and avoid comparing systems tuned with different $K$ or $n$ on reported scores. Methods that look attractive can fail: empirical-Bayes shrinkage of disclosed scores assumes independent noise, and the Ladder-style report depends on a free step size.

\textbf{Limitations.} (i) The real-LLM evidence is one 1.5B-parameter model on a pool of 2000 MMLU-Pro questions on CPU; no larger model, second benchmark or second task was run, so we make no statement about frontier systems. (ii) The candidates are prompt formats and two-prompt ensembles; their true spread is small and they are correlated, which is a low-signal regime; real prompt/program search with instruction proposals may have larger spread and adaptive proposals that our three procedures (exhaustive, halving, greedy) only approximate. (iii) T1 is a replay on cached outcomes, \emph{not} a live run of DSPy/MIPRO; the live-tool tier of our plan was not done. (iv) The pool is fixed, and the uncertainty in T1 is over search sets given the pool; between-pool uncertainty was not estimated; the block SEs reflect variation over search sets given the pool, not the pool. (v) T0 is a model of the mechanism, not calibrated to T1 (its gap is larger than the real one in 28 of 30 nearest cells, median ratio 1.6). (vi) The extreme-value formula is heuristic. (vii) Selection-aware inference (M4) is implemented only for exhaustive search; for halving and greedy only M0--M3 and M5 are available. (viii) Coverage and decision rates in cells with few replicates in a regime (e.g.\ false-claim regimes at large $\tau$) have wide intervals (Table~\ref{tab:D}). (ix) Compute (Table~\ref{tab:F}, appendix): the synthetic simulations took 3{,}807 s of summed job wall time, scoring 24 prompts on 2000 items (48{,}000 model calls) took 81{,}172 s elapsed on 8 parallel 2-vCPU CPU jobs, and the replays on cached outcomes 1{,}217 s; the 0.5B model arm, a second model and a live DSPy/MIPRO run were not done.
